\scanpage{60}
\uline{Corollaire 3}\quad Si $H$ est comme précédemment,
\unsure{on a} $\underline{\operatorname{Ext}}^1(H,\underline W)=0$
$(\underline W=\varinjlim_n W_n)$.
\figure{\begin{tikzpicture}\draw (0,0)--(0,-.5);\end{tikzpicture}}

\margin{\uline{Question} Si $S$ est affine, a-t-on
$\operatorname{Ext}^1(H,\underline W)=0$
i.e. $H^1(S,\underline{\operatorname{Hom}}(H,\underline W))=0$ ??
\uuline{Non} déjà pour $H=\alpha_p$, si $S=\operatorname{Spec}k$,
$k$ corps non parfait.}

\uline{Proposition 4} Soit $G$ un groupe algébrique
\deleted{\unsure{fini}} \inserted{fini, radiciel et unipotent}
sur un corps $k$ de car. $p>0$ \inserted{(pas néc. parfait)}
tel que $G$ soit annulé par $V^n$. Alors

\deleted{a) $\underline{\operatorname{Hom}}_{\mathrm{gp}}(G,W_n)$
est un schéma en groupes lisse \inserted{connexe}
\deleted{sur $k$ et dimension \unsure{$r=\dim G/G^0$}}
(et $\overline{\mathrm m}$ isomorphe à $\unsure{\mathbb E_k^r|\mathbb A_k^r}$
pour $r$ \deleted{$\unsure{\dim_k G}$}
$\operatorname{rg}G=p^r$)}

b) $\underline{\operatorname{Ext}}^1_{\unsure{\mathrm{gp}}}
(G,\underrightarrow W)=0$

Le b) provient du cor. 3, compte tenu que $G$ se dévisse en
groupes isom. à $\alpha_p$.

\deleted{Pour prouver a), on procède par récurrence sur la longueur
$r$ d'une \inserted{telle} suite de composition. On a donc
\[
0-G'\longrightarrow G-\alpha_p-0
\]
d'où une suite exacte
\[
0-\underline{\operatorname{Hom}}(\alpha_p,\underrightarrow W)
-\underline{\operatorname{Hom}}(G,\underline W)
-\underline{\operatorname{Hom}}(G',\underline W)\longrightarrow0
\]
compte tenu que $\underline{\operatorname{Ext}}^1(\alpha_p,\underline W)=0$.
Or les trois termes ne changent pas si on y remplace $\underline W$
par $W_n$. Enfin, on applique l'hyp. de récurrence à
$\underline{\operatorname{Hom}}(G',W_n)$, et on note que
$\underline{\operatorname{Hom}}(\alpha_p,W_n)\simeq G_a$.
On gagne!

5. Soit $\mathbb D(G)=\underline{\operatorname{Hom}}(G,\underrightarrow W)$.
C'est un $W_\infty$-Module, muni des endom $V$ et $F$ satisfaisant
les conditions habituelles (provenant des relations entre $V,F$
et les $\lambda\in W_\infty(S)$, qu'on a dans $\underrightarrow W$)
\[
(5.1)\qquad F\lambda=\lambda^{(p)}F,\quad
\lambda V=V\lambda^p,\qquad FV=VF=p.\operatorname{id}
\]
}
\margin{70}
